Das Auseinandersetzen mit Transistoren ist unabdingbar, um ein Gefühl für die Hintergründe und Funktionsweisen der technischen Bauteile zu bekommen, deren Verfeinerung, Verbesserung und Weiterentwicklung die Digitale Revolution in einer Dimension ermöglicht, die alle Bereiche der menschlichen Zivilisation von Grund auf verändert. 

Zudem sind Löt- und Grundlagenkenntnisse elektrischer Schaltungen ein wesentlicher Bestandteil der Realisation von Experimenten und eigenen Anwendungen, das Beschäftigen mit diesen Bereichen ist also von großer Bedeutung für Physiker:innen (und natürlich auch von Interesse für von Elektronik faszinierte Menschen). Außerdem wurde weiter Erfahrung im Umgang mit Oszilloskopen gesammelt, was auch ein wichtiger Skill für Physiker:innen ist.

Final liegt der Charme dieses Versuches in dem sukzessiven Verstehen der Bestandteile und Aufbauen eines hochgradig fähigen Systems, dem PID-Regler, dessen Schaltplan \cite[vgl.][11]{Skript_01} zwar auf den ersten Blick überwältigend aussieht, auf den zweiten Blick jedoch nur aus geschickter Kombination bekannter Bauteilen besteht. Ein PID-Regler basiert auf der eleganten Nutzung von Operationsverstärkern, deren Innenleben zwar nur aus Widerständen, Kondensatoren und Transistoren besteht, jedoch alles andere als trivial ist:
\xfigure{htbp}{width=0.8\textwidth}{OpAmpTransistor.png}{\emph{Meme:} Which part don't you understand?\autocite{OpAmpTransistorLevelColored}}{ist der Operationsverstärker wirklich bekannt? \autocite{OpAmpTransistorLevelColored}}

\paragraph{PID-Grundlagen:} Es gibt gute Youtube-Videos über PID-Controlling: kurzer Überblick: \href{https://youtu.be/wkfEZmsQqiA?si=atU3_SOSYyGcCa4w}{\emph{What Is PID Control?}} mit Beispiel und hier ein etwas komplexeres Beispiel \href{https://www.youtube.com/watch?v=XfAt6hNV8XM}{\emph{Simple Examples of PID Control}}, das die Eigenschaften von P-, I-, D-, PID-Controllern aufzeigt.